\documentclass[aspectratio=169]{beamer}
\input{header}
\coursecode{JKSSB}

\title{Current Cybersecurity, Cryptography \& Cyber Law}
\subtitle{Lesson 64 --- Crypto, Firewalls and Law}
\author{JKSSB Computer Awareness}
\date{\today}

\begin{document}
\frame{\titlepage}

\begin{frame}{Where This Lesson Fits}
  \begin{itemize}
    \item The syllabus also tests \key{computer virus \& latest anti-virus},
      \key{terms and abbreviations used in IT}, and the \key{role of IT in governance}.
    \item This lesson covers the current-affairs edge of that scope: threats,
      crypto basics, firewalls, and cyber law at \key{awareness level}.
    \item Depth stays at recall and classification (\muted{L1/L2}), with one
      short scenario at most.
  \end{itemize}
  \begin{block}{The one idea}
    Know what each threat, tool, and term \key{does} --- and what its nearest
    neighbour does instead.
  \end{block}
\end{frame}

\begin{frame}{Malware Family: Virus First}
  \begin{itemize}
    \item \key{Malware} is hostile software: virus, worm, Trojan horse,
      ransomware, spyware.
    \item A \key{virus} needs a host program and human action to spread; a
      \key{worm} spreads on its own across a network.
    \item A \key{Trojan horse} looks like useful software but hides a harmful
      payload; it does \key{not} self-replicate.
  \end{itemize}
  \begin{block}{Keep the pair straight}
    Virus needs a host and help to spread; worm spreads by itself.
  \end{block}
\end{frame}

\begin{frame}{Parasitic (File-Infector) Virus}
  \begin{itemize}
    \item A \key{parasitic virus}, also called a \key{file-infector virus},
      attaches itself to an executable file.
    \item When the infected program runs, the virus code runs \key{first},
      then hands control back to the host program.
    \item It spreads as infected files are \key{shared or copied} to other
      machines.
    \item An \key{anti-virus} detects, blocks, and removes such infections;
      its database must be kept \key{updated}.
  \end{itemize}
  \begin{alertblock}{Trap alert}
    A file infector attaches to files. It is not a worm and not a Trojan.
  \end{alertblock}
\end{frame}

\begin{frame}{Four Threats Side by Side}
  \begin{center}
    \begin{tabular}{p{0.18\textwidth}p{0.62\textwidth}}
      \toprule
      \textbf{Threat} & \textbf{What happens} \\
      \midrule
      Phishing & Fraudulent mail or message tricks the user into revealing passwords or OTPs \\
      Ransomware & Malicious code locks or encrypts files, then demands payment \\
      DDoS & Flood of traffic from many machines overwhelms a server \\
      Hacking & Unauthorised access to a computer or account \\
      \bottomrule
    \end{tabular}
  \end{center}
  \vspace{0.25cm}
  \muted{Each row is a different mechanism; the exam swaps their definitions.}
\end{frame}

\begin{frame}{Phishing vs Ransomware vs DDoS vs Hacking}
  \begin{itemize}
    \item \key{Phishing} attacks the \key{user}: fake bank or KYC mail, fake
      login page, stolen credentials.
    \item \key{Ransomware} attacks the \key{data}: files are encrypted and a
      ransom is demanded for the key.
    \item \key{DDoS} (Distributed Denial of Service) attacks
      \key{availability}: the service is flooded until genuine users cannot
      reach it.
    \item \key{Hacking} is the broad term for \key{unauthorised access}; the
      other three are specific methods or goals.
  \end{itemize}
  \begin{exampleblock}{Office desktop example}
    A clerk gets a fake ``verify your account'' mail (phishing), while the
    office server slows to a halt under flood traffic (DDoS) --- different
    threats, different defences.
  \end{exampleblock}
\end{frame}

\begin{frame}{Cryptography in One Frame}
  \begin{itemize}
    \item \key{Cryptography} protects information by transforming it so only
      the intended recipient can read it.
    \item \key{Plaintext} is the readable message; \key{ciphertext} is the
      transformed message; a \key{key} controls the transformation.
    \item \key{Encryption} converts plaintext to ciphertext;
      \key{decryption} converts it back.
    \item Two broad families: \key{symmetric} (same key both sides) and
      \key{asymmetric} (public key and private key pair).
  \end{itemize}
\end{frame}

\begin{frame}{DES: The 64-Bit Block Fact}
  \begin{itemize}
    \item \key{DES} stands for Data Encryption Standard.
    \item It is a \key{symmetric} cipher: the same key encrypts and decrypts.
    \item DES works on \key{64-bit blocks} of data at a time.
    \item Its key is 64 bits long, of which \key{56 bits} are effective (8
      bits are parity).
  \end{itemize}
  \begin{alertblock}{Trap alert}
    DES block size is \key{64 bits}. Do not confuse block size with the 56-bit
    effective key.
  \end{alertblock}
\end{frame}

\begin{frame}{Encryption vs Digital Signature (TRAP-15)}
  \begin{itemize}
    \item \key{Encryption} provides \key{confidentiality}: only the intended
      recipient can read the message.
    \item A \key{digital signature} provides \key{authenticity and integrity}:
      it proves who sent the message and that it was not altered.
    \item A signature does \key{not} hide the message content; encryption
      does \key{not} by itself prove the sender's identity.
  \end{itemize}
  \begin{alertblock}{Trap alert}
    Digital signature $\ne$ encryption. Signature = who sent it and unchanged;
    encryption = who can read it. (TRAP-15/TRM-11)
  \end{alertblock}
\end{frame}

\begin{frame}{Firewalls: What They Do}
  \begin{itemize}
    \item A \key{firewall} is a barrier between a trusted internal network and
      an untrusted external network.
    \item It permits or denies traffic according to a \key{rule set}.
    \item It filters \key{packets and connections}, not file-infecting viruses
      inside allowed downloads --- that is the anti-virus's job.
    \item Three types tested: packet-filtering, stateful inspection, and
      application (proxy) firewalls.
  \end{itemize}
\end{frame}

\begin{frame}{Three Firewall Types}
  \begin{center}
    \begin{tabular}{p{0.26\textwidth}p{0.28\textwidth}p{0.32\textwidth}}
      \toprule
      \textbf{Type} & \textbf{Checks} & \textbf{Level} \\
      \midrule
      Packet-filtering & Source/destination address, port, protocol & Network layer, per packet \\
      Stateful inspection & Same headers \key{plus connection state} & Tracks sessions, not single packets \\
      Application / proxy & Content and commands of the application & Application layer; acts as go-between \\
      \bottomrule
    \end{tabular}
  \end{center}
  \vspace{0.25cm}
  \muted{Order of depth: packet header only, then header plus state, then full
  application content.}
\end{frame}

\begin{frame}{Firewall Ordering Drill}
  \begin{itemize}
    \item \key{Packet-filtering}: fastest and simplest; decides packet by
      packet; cannot judge a session as a whole.
    \item \key{Stateful inspection}: remembers open connections; drops packets
      that do not belong to a known session.
    \item \key{Application/proxy firewall}: sits between client and server and
      inspects application commands (for example web or mail requests).
    \item Exam shorthand: header only $\rightarrow$ header plus state
      $\rightarrow$ full content via proxy.
  \end{itemize}
  \begin{block}{Keep the triple straight}
    Filtering checks the envelope; stateful checks the conversation; proxy
    reads the letter inside.
  \end{block}
\end{frame}

\begin{frame}{Internet Governance Concerns}
  \begin{itemize}
    \item The Internet has \key{no single owner}; addressing and names are
      coordinated internationally (for example by ICANN for names and numbers).
    \item Governance concerns tested at awareness level: \key{security} of
      networks, \key{privacy} of personal data, fair \key{access}, and which
      country's \key{law} applies to cross-border acts.
    \item The \key{role of IT in governance} (e-governance) is the positive
      side: online services, transparency, and delivery to citizens.
  \end{itemize}
\end{frame}

\begin{frame}{Privacy and Data Protection}
  \begin{itemize}
    \item \key{Personal data} is any data that identifies a person: name with
      phone number, biometrics, account details.
    \item Protection principles at awareness level: collect with
      \key{consent}, use only for the stated \key{purpose}, keep data
      \key{secure}, and retain it no longer than needed.
    \item \key{Privacy settings}, strong passwords, and OTP discipline are the
      user-side defences the exam names.
    \item Leaking or misusing personal data can also attract \key{legal}
      consequences under cyber law.
  \end{itemize}
\end{frame}

\begin{frame}{Copyright Act: Awareness-Level Map}
  \begin{itemize}
    \item The \key{Copyright Act} protects original literary, artistic,
      musical, and software works from copying without permission.
    \item Awareness-level mapping only: \key{meaning of copyright} (what rights
      the owner holds), \key{first ownership} (who owns it first),
      \key{infringement} (what counts as violation), \key{fair dealing}
      (limited permitted uses such as study), \key{offences} (punishment for
      wilful infringement).
    \item Software piracy --- copying licensed software without permission ---
      is the exam's standard infringement example.
  \end{itemize}
  \begin{block}{Exam depth}
    Match the \key{concept} to its section area; section numbers are recall
    only where the question itself supplies them.
  \end{block}
\end{frame}

\begin{frame}{Abbreviation and Term Ladder}
  \begin{itemize}
    \item \key{DES} --- Data Encryption Standard; 64-bit blocks.
    \item \key{DDoS} --- Distributed Denial of Service; flood from many machines.
    \item \key{OTP} --- One Time Password; phishing target.
    \item \key{ICANN} --- Internet Corporation for Assigned Names and Numbers.
    \item \key{IP} --- Internet Protocol; the address a packet filter reads.
  \end{itemize}
  \begin{alertblock}{Trap alert}
    DDoS needs many machines (distributed); a single-machine flood is only DoS.
  \end{alertblock}
\end{frame}

\begin{frame}{Trap Alert: Facts to Keep Straight}
  \begin{itemize}
    \item DES block size is \key{64 bits} (56-bit effective key) --- not 128.
    \item Digital signature $\ne$ encryption: authenticity/integrity vs
      confidentiality. (TRAP-15)
    \item File infector attaches to \key{executable files}; it is not a worm.
    \item Packet filter checks headers only; stateful adds \key{connection
      state}; proxy inspects \key{application content}.
    \item Phishing steals \key{credentials}; ransomware locks \key{files};
      DDoS blocks \key{access}; hacking is unauthorised access generally.
  \end{itemize}
\end{frame}

\begin{frame}{Lesson 64: Recall}
  \begin{itemize}
    \item A parasitic (file-infector) virus attaches to executable files and
      runs when the host program runs.
    \item DES (Data Encryption Standard) is symmetric and works on 64-bit blocks.
    \item Firewalls: packet-filtering (headers), stateful (plus connection
      state), application/proxy (application content).
    \item Phishing tricks the user; ransomware encrypts files for ransom; DDoS
      floods a service; hacking is unauthorised access.
    \item Encryption gives confidentiality; a digital signature gives
      authenticity and integrity.
    \item Governance concerns: security, privacy, access, jurisdiction;
      copyright protects owners against infringement, with fair-dealing exceptions.
  \end{itemize}
\end{frame}
\end{document}
